\documentclass[12pt,a4paper]{article}

\usepackage[a4paper,top=2.2cm,bottom=2.3cm,left=2.2cm,right=2.2cm]{geometry}
\usepackage{graphicx}
\usepackage{booktabs}
\usepackage{longtable}
\usepackage{array}
\usepackage[backend=biber,style=numeric]{biblatex}
\addbibresource{references.bib}
\usepackage{amsmath,amssymb}
\usepackage{float}
\usepackage{xcolor}
\usepackage{pgfplots}
\usepackage{caption}
\usepackage{enumitem}
\usepackage{fancyhdr}
\usepackage{caption}
\usepackage{subcaption}
\usepackage{tikz}
\usetikzlibrary{arrows.meta,positioning,shapes.geometric,calc,fit}
\usepackage{hyperref}
\hypersetup{
  colorlinks=true,
  linkcolor=blue!45!black,
  urlcolor=blue!55!black,
  citecolor=green!35!black,
  pdftitle={گزارش توجیه داخلی عکس هوایی},
  pdfsubject={فتوگرامتری و آشکارسازی نیمه‌خودکار نشانه‌های فیدوشال}
}
\geometry{margin=2.5cm}
\pgfplotsset{compat=1.18}
\usepackage{xepersian}


\settextfont{XB Niloofar}
\setlatintextfont{Times New Roman}
\setdigitfont{Yas}


\pagestyle{fancy}
\fancyhf{}
\fancyhead[R]{\small گزارش پروژه امتیازی کابردهای فتوگرامتری}
\fancyhead[L]{\small نیم‌سال دوم ۱۴۰5–۱۴۰4}
\fancyfoot[C]{\thepage}
\fancyfoot[L]{دانشکده مهندسی نقشه برداری}
\fancyfoot[R]{دانشگاه صنعتی خواجه نصیر‌الدین طوسی}
\renewcommand{\footrulewidth}{0.4pt}
\renewcommand{\arraystretch}{1.35}
\setlength{\parindent}{0pt}
\setlength{\parskip}{0.65em}
\captionsetup{font=small,labelfont=bf}

\definecolor{projectblue}{RGB}{39,92,142}
\definecolor{projectgreen}{RGB}{46,125,92}
\definecolor{projectorange}{RGB}{202,113,42}
\definecolor{projectgray}{RGB}{242,245,247}

\newcommand{\code}[1]{\lr{\texttt{\detokenize{#1}}}}
\newcommand{\unit}[1]{\text{\lr{#1}}}
\renewcommand{\unit}[1]{\ensuremath{\mathrm{#1}}}



\begin{document}
\begin{titlepage}
\centering
\begin{figure}[H]
\centering

\includegraphics[width=0.4\textwidth]{KNTU.png}
\end{figure}
{\large دانشکده‌ی مهندسی نقشه‌برداری (ژئودزی و ژئوماتیک)}\\[2cm]

\rule{\textwidth}{1.5pt}\\[0.4cm]
{\LARGE\bfseries (پروژه امتیازی) گزارش فرایند توجیه داخلی یک عکس هوایی آنالوگ }\\[0.2cm]
{\LARGE \textbf{}}\\[0.7cm]

{\large  آشکارسازی نیمه‌خودکار فیدوشال مارکها و ارزیابی تبدیل‌های هندسی}\\[2cm]
\rule{\textwidth}{1.5pt}\\[1cm]
{\Large کاربرد‌های فتوگرامتری}\\[0.33cm]
{\large استاد: جناب دکتر عبادی }\\[0.3cm]
{\large استادیار: خانم مهندس نعنجم }\\[0.7cm]
{\large دانشجو: علی بهروز(40215613) }\\[0.8cm]

{\large نیم‌سال دوم تحصیلی ۱۴۰5– ۱۴۰4 }\\[3cm]



\vfill
\end{titlepage}
\newpage
\thispagestyle{empty}
\begin{center}
    \vspace*{\fill}
    
\includegraphics[width=0.9\textwidth]{besm-nagsh.ir-154.jpg}
    \vspace*{\fill}
\end{center}
\newpage
\thispagestyle{empty}

\begin{abstract}
در این پروژه، توجیه داخلی یک عکس هوایی اسکن‌شده از شهر همدان با استفاده از اطلاعات
کالیبراسیون دوربین \lr{Wild RC20} و هشت فیدوشال مارک انجام شده است. برنامه
\code{interior_orientation.py} ابتدا دو نشانه\lr{(Mark)} را به‌صورت دستی دریافت و مکان آن‌ها را
با همبستگی متقابل نرمال‌شده\lr{(Normalized Cross Correlation)}
\cite{lewis1995fast,opencv_template_matching} پالایش می‌کند. سپس با حل یک تبدیل تشابهی دوبعدی، موقعیت
تقریبی سایر نشانه‌ها پیش‌بینی می‌شود و مختصات دقیق آن‌ها با تطبیق الگو و برآورد
زیرپیکسلی تعیین می‌گردد. در ادامه، تبدیل‌های افاین و پروجکتیو میان دستگاه مختصات
پیکسلی و دستگاه مختصات کالیبره دوربین با روش کمترین مربعات برآورد و مقایسه می‌شوند.
در نهایت، مرکز فیدوشال از تقاطع دو خط واصل نشانه‌های روبه‌رو محاسبه شده است. خطای
\lr{RMSE} برابر با \(0.00972\,\mathrm{mm}\) برای تبدیل افاین و
\(0.00936\,\mathrm{mm}\) برای تبدیل پروجکتیو به‌دست آمد که معادل حدود نیم پیکسل است.
\end{abstract}

\textbf{واژگان کلیدی:}
توجیه داخلی، فتوگرامتری هوایی، فیدوشال مارک، همبستگی متقابل نرمال‌شده،
تبدیل افاین، هموگرافی، نقطه اصلی.

\newpage
\clearpage
\tableofcontents
\clearpage
\newpage

\section{مقدمه}

توجیه داخلی فرایندی است که رابطه هندسی میان دستگاه مختصات تصویر دیجیتال
\((r,c)\) و دستگاه مختصات متریک دوربین \((X,Y)\) را بازسازی می‌کند. در دوربین‌های
آنالوگ، مختصات فیدوشال‌ها در گواهی کالیبراسیون معلوم است؛ اما پس از اسکن فیلم،
دوران، انتقال، تغییر مقیاس، ناهمسانی مقیاس در دو راستا و اعوجاج‌های کوچک اسکنر باعث
می‌شوند مختصات پیکسلی مستقیماً قابل استفاده نباشد. بنابراین باید فیدوشال مارک‌ها
در تصویر شناسایی و یک تبدیل هندسی میان دو دستگاه مختصات برآورد شود
\cite{wolf2014elements}.

هدف این پروژه، پیاده‌سازی یک گردش‌کار نیمه‌خودکار برای انجام این فرایند است. داده‌های
اصلی شامل تصویر \code{1614.tif}، پرونده کالیبراسیون \code{RC20.txt}، فراداده پرواز
\code{meta.txt} و خروجی عددی \code{results.txt} هستند. شکل
\ref{fig:detected-image} نتیجه نهایی آشکارسازی نشانه‌ها و مرکز فیدوشال را نمایش
می‌دهد.

\section{داده‌ها و مشخصات پروژه}

\subsection{مشخصات تصویر و پرواز}

تصویر ورودی یک تصویر خاکستری هشت‌بیتی با ابعاد
\(12173\times12203\) پیکسل است. اطلاعات پرواز و تصویر مطابق
\code{meta.txt} در جدول \ref{tab:flight-data} خلاصه شده‌اند.

\begin{table}[H]
\centering
\caption{مشخصات تصویر و پرواز}
\label{tab:flight-data}
\begin{tabular}{>{\raggedleft\arraybackslash}p{0.42\textwidth}
                >{\centering\arraybackslash}p{0.38\textwidth}}
\toprule
\textbf{پارامتر} & \textbf{مقدار} \\
\midrule
منطقه تصویربرداری & شهر همدان \\
مقیاس اسمی عکس & \(1:5000\) \\
فاصله کانونی & \(153.24\,\mathrm{mm}\) \\
ارتفاع پرواز از سطح متوسط منطقه، \(H'\) & \(766.200\,\mathrm{m}\) \\
ارتفاع متوسط منطقه & \(1883.73\,\mathrm{m}\) \\
باز عکسی متوسط، \(b\) & \(77.5\,\mathrm{mm}\) \\
ابعاد تصویر رقومی & \(12173 \times 12203\,\unit{px}\) \\
\bottomrule
\end{tabular}
\end{table}


رابطه مقیاس عکس به‌صورت زیر است:
\begin{equation}
  m=\frac{f}{H'}=\frac{0.15324}{766.2}\simeq\frac{1}{5000}.
\end{equation}
همچنین، باز هوایی تقریبی از رابطه \(s=b/m\) به‌دست می‌آید:
\begin{equation}
  s=77.5\,\mathrm{mm}\times5000=387.5\,\mathrm{m},
  \qquad
  \frac{B}{H'}\simeq0.506.
\end{equation}

\subsection{مشخصات دوربین}

دوربین از نوع فیلمی \lr{RC20} با شماره سریال \(13180\) است. اطلاعات مهم پرونده
\code{RC20.txt} عبارت‌اند از:
\begin{itemize}[rightmargin=1.2em]
  \item فاصله کانونی کالیبره: \(153.240\,\mathrm{mm}\)؛
  \item مختصات نقطه اصلی تقارن نسبت به مرکز فیدوشال:
  \((x_{\mathrm{PPS}},y_{\mathrm{PPS}})=(-0.004,\,+0.005)\,\mathrm{mm}\)؛
  \item مدل اعوجاج ثبت‌شده: اعوجاج شعاعی ساده؛
  \item تعداد  فیدوشال مارک‌ها: هشت نشانه.
\end{itemize}

\begin{table}[H]
\centering
\caption{مختصات کالیبره فیدوشال مارک‌ها در دستگاه دوربین}
\label{tab:fiducial-calibration}
\begin{tabular}{ccc@{\qquad}ccc}
\toprule
\textbf{شناسه} & \(X~[\unit{mm}]\) & \(Y~[\unit{mm}]\) &
\textbf{شناسه} & \(X~[\unit{mm}]\) & \(Y~[\unit{mm}]\) \\
\midrule
1 & \(106.001\) & \(-106.001\) & 5 & \(0.003\) & \(-110.000\) \\
2 & \(-106.001\) & \(-106.001\) & 6 & \(-109.997\) & \(-0.001\) \\
3 & \(-106.002\) & \(106.003\) & 7 & \(0.003\) & \(110.001\) \\
4 & \(106.005\) & \(106.006\) & 8 & \(110.006\) & \(0.000\) \\
\bottomrule
\end{tabular}
\end{table}

\begin{figure}[H]
\centering
\begin{tikzpicture}[
    scale=0.92,>=Latex,
    fmark/.style={projectblue,very thick},          % fiducial marks
    fnum/.style={fill=white,draw=projectblue!40,     % number badges
                 rounded corners=1pt,inner sep=1.6pt,
                 font=\footnotesize\bfseries},
    axisstyle/.style={thick,gray!60}                 % faint axis grid
]

  % ---- Frame with soft shadow ----
  \fill[black!25,rounded corners=3pt] (-5.25,-5.55) rectangle (5.55,5.25);
  \draw[fill=projectgray,rounded corners=3pt,draw=black!30]
        (-5.4,-5.4) rectangle (5.4,5.4);

  % ---- Faint interior grid for reference ----
  \foreach \g in {-4,-3,-2,-1,1,2,3,4}{
    \draw[axisstyle,opacity=0.15] (\g,-5.1) -- (\g,5.1);
    \draw[axisstyle,opacity=0.15] (-5.1,\g) -- (5.1,\g);
  }

  % ---- Diagonal lines joining the corner fiducials (1-3, 2-4) ----
  \draw[projectorange!70,thick,dash pattern=on 4pt off 2pt]
        (4.1,-4.1) -- (-4.1,4.1);
  \draw[projectorange!70,thick,dash pattern=on 4pt off 2pt]
        (-4.1,-4.1) -- (4.1,4.1);

  % ---- Axes joining the edge-fiducial pairs (5-7, 6-8) ----
  \draw[projectorange,very thick,dashed] (-4.55,0) -- (4.55,0);
  \draw[projectgreen,very thick,dashed] (0,-4.55) -- (0,4.55);

  % ---- Coordinate axes on top ----
  \draw[->,thick,black!80] (-4.7,0) -- (5.0,0) node[below right] {\(X\)};
  \draw[->,thick,black!80] (0,-4.7) -- (0,5.0) node[above left] {\(Y\)};

  % ---- Fiducial marks: crosshair + circle (more realistic look) ----
  \foreach \name/\x/\y in {
    1/4.1/-4.1, 2/-4.1/-4.1, 3/-4.1/4.1, 4/4.1/4.1,
    5/0/-4.55,  6/-4.55/0,   7/0/4.55,   8/4.55/0}
  {
    \draw[fmark] (\x,\y) circle (0.16);
    \draw[fmark] (\x-0.3,\y) -- (\x+0.3,\y);   % horizontal hairline
    \draw[fmark] (\x,\y-0.3) -- (\x,\y+0.3);   % vertical hairline
    \fill[projectblue] (\x,\y) circle (0.04);  % center dot
    \node[fnum] at ({1.13*\x},{1.13*\y}) {\name};
  }

  % ---- Principal / fiducial centre ----
  \draw[red,very thick] (-0.22,0) -- (0.22,0);
  \draw[red,very thick] (0,-0.22) -- (0,0.22);
  \fill[red] (0,0) circle (2.4pt);
  \node[fill=white,draw=red!40,rounded corners=1pt,inner sep=2pt,
        font=\footnotesize,above right=2pt and 2pt] at (0,0) {\lr{مرکز فیدوشال}};

\end{tikzpicture}
\caption{آرایش اسمی هشت فیدوشال مارک و زوج‌های روبه‌رو \(6\!-\!8\) و \(5\!-\!7\)}
\label{fig:fiducial-layout}
\end{figure}


\section{گردش‌کار محاسباتی}

شکل \ref{fig:workflow} توالی کلی پردازش در برنامه را نشان می‌دهد. هسته روش شامل
شش مرحله اصلی است.

\begin{figure}[H]
\centering
\begin{tikzpicture}[
  node distance=7mm and 10mm,
  box/.style={
    rectangle,rounded corners=3pt,draw=projectblue,very thick,
    fill=blue!4,text width=0.82\textwidth,minimum height=10mm,
    align=center,inner sep=5pt
  },
  arrow/.style={-{Latex[length=3mm]},very thick,projectblue}
]
  \node[box] (input) {\rl{خواندن تصویر خاکستری و مختصات کالیبره فیدوشال‌ها}};
  \node[box,below=of input] (manual) {\rl{اندازه‌گیری دستی دو فیدوشال و پالایش موضعی آن‌ها با \lr{NCC}}};
  \node[box,below=of manual] (similarity) {\rl{حل تبدیل تشابهی چهارپارامتری و پیش‌بینی مکان هر هشت نشانه}};
  \node[box,below=of similarity] (matching) {\rl{ساخت الگوی حلقوی، جست‌وجوی موضعی و آشکارسازی زیرپیکسلی}};
  \node[box,below=of matching] (models) {\rl{برآورد کمترین مربعات تبدیل افاین و هموگرافی پروجکتیو}};
  \node[box,below=of models] (pp) {\rl{محاسبه تقاطع خطوط نشانه‌های روبه‌رو و تعیین مرکز فیدوشال}};
  \node[box,below=of pp] (output) {\rl{گزارش ماتریس‌ها، \lr{RMSE}، امتیازهای \lr{NCC} و ترسیم نتیجه}};

  \draw[arrow] (input) -- (manual);
  \draw[arrow] (manual) -- (similarity);
  \draw[arrow] (similarity) -- (matching);
  \draw[arrow] (matching) -- (models);
  \draw[arrow] (models) -- (pp);
  \draw[arrow] (pp) -- (output);
\end{tikzpicture}
\caption{نمودار گردش‌کار برنامه \code{interior_orientation.py}}
\label{fig:workflow}
\end{figure}

\subsection{اجرای بازتولیدپذیر برنامه}

برنامه به کتابخانه‌های \lr{NumPy}، \lr{OpenCV} و \lr{Matplotlib} نیاز دارد
\cite{harris2020numpy,bradski2000opencv,hunter2007matplotlib}. یک
نمونه اجرای تعاملی با انتخاب فیدوشال‌های ۱ و ۳ به‌صورت زیر است:
\begin{LTR}
\small
\begin{verbatim}
python3 interior_orientation.py 1614.tif RC20.txt --click-ids 1 3 \
  --save-figure Figure_1.png
\end{verbatim}
\end{LTR}
در صورت معلوم بودن مختصات اولیه، می‌توان به‌جای انتخاب گرافیکی، گزینه
\code{--manual ID ROW COLUMN} را دقیقاً دو بار به فرمان افزود. اندازه الگو و شعاع
جست‌وجو نیز به‌ترتیب با \code{--template-size} و \code{--search-radius} قابل تغییر
هستند.

\subsection{مرحله اول: دریافت و پالایش دو فیدوشال اولیه}

کاربر باید دقیقاً دو فیدوشال را یا با گزینه \code{--manual} وارد کند یا با گزینه
\code{--click-ids} در محیط تعاملی انتخاب نماید. در حالت تعاملی، ابتدا محل تقریبی
فیدوشال روی نمای کلی و سپس مرکز دقیق آن در پنجره بزرگ‌نمایی انتخاب می‌شود. برنامه برای
پالایش این نقاط نیز همان سازوکار تطبیق الگو را به‌کار می‌گیرد.

\begin{figure}[H]
\centering
\includegraphics[width=0.97\textwidth]{one.png}
\caption{مرحله ابتدایی و صفحه اولیه تعیین فیدوشال مارک‌های اول و سوم به صورت دستی/تعاملی}
\label{fig:first}
\end{figure}

\begin{figure}[H]
\centering
\includegraphics[width=0.97\textwidth]{two.png}
\caption{انتخاب و تعیین موقعیت پیکسلی فیدوشال مارک اول}
\label{fig:second}
\end{figure}

\begin{figure}[H]
\centering
\includegraphics[width=0.97\textwidth]{three.png}
\caption{انتخاب و تعیین موقعیت پیکسلی فیدوشال مارک سوم}
\label{fig:third}
\end{figure}

\subsection{مرحله دوم: تبدیل تشابهی و پیش‌بینی اولیه}

برای دو نقطه متناظر، تبدیل تشابهی زیر میان مختصات پیکسلی \((r,c)\) و مختصات دوربین
\((X,Y)\) حل می‌شود:
\begin{equation}
\begin{aligned}
X &= a\,r-b\,c+t_x,\\
Y &= b\,r+a\,c+t_y.
\end{aligned}
\label{eq:conformal}
\end{equation}
چهار پارامتر \(a,b,t_x,t_y\) با حل مستقیم چهار معادله حاصل از دو نقطه تعیین
می‌شوند. سپس تبدیل معکوس برای پیش‌بینی مکان همه فیدوشال‌ها استفاده می‌شود:
\begin{equation}
\begin{aligned}
r &=
\frac{a(X-t_x)+b(Y-t_y)}{a^2+b^2},\\
c &=
\frac{-b(X-t_x)+a(Y-t_y)}{a^2+b^2}.
\end{aligned}
\end{equation}
این تبدیل صرفاً مکان تقریبی جست‌وجو را فراهم می‌کند و مدل نهایی توجیه داخلی نیست
\cite{wolf2014elements}.

\subsection{مرحله سوم: ساخت الگو و پنجره جست‌وجو}

برای هر فیدوشال، یک الگوی حلقوی مصنوعی با اندازه پیش‌فرض \(101\times101\) پیکسل ساخته
می‌شود. شعاع حلقه تقریباً \(0.4\) عرض الگو است و برای کاهش حساسیت به نویز، فیلتر
گاوسی روی آن اعمال می‌شود. همچنین امکان معرفی الگوی مستقل برای هر شناسه با گزینه
\code{--external-template} وجود دارد. پنجره جست‌وجو در پیرامون محل پیش‌بینی‌شده و
با شعاع پیش‌فرض \(50\) پیکسل تشکیل می‌شود.

\subsection{مرحله چهارم: تطبیق الگو و مختصات زیرپیکسلی}

معیار تطبیق، ضریب همبستگی متقابل نرمال‌شده از نوع
\code{TM_CCOEFF_NORMED} است
\cite{lewis1995fast,opencv_template_matching}. به‌طور مفهومی، این معیار را می‌توان به‌شکل
\begin{equation}
\operatorname{NCC}(u,v)=
\frac{\sum\limits_{i,j}
\left[I_{u,v}(i,j)-\overline{I}_{u,v}\right]
\left[T(i,j)-\overline{T}\right]}
{\sqrt{\sum\limits_{i,j}\left[I_{u,v}(i,j)-\overline{I}_{u,v}\right]^2}
 \sqrt{\sum\limits_{i,j}\left[T(i,j)-\overline{T}\right]^2}}
\end{equation}
نوشت. بیشینه سطح همبستگی به‌عنوان محل فیدوشال پذیرفته می‌شود، مشروط بر آن‌که امتیاز
از آستانه پیش‌فرض \(0.35\) بیشتر باشد و بیشینه در مرز پنجره قرار نگیرد.

برای افزایش دقت، یک سهمی یک‌بعدی روی سه مقدار مجاور قله در هر راستا برازش می‌شود:
\begin{equation}
\delta=-\frac{1}{2}
\frac{q_{+1}-q_{-1}}{q_{+1}-2q_0+q_{-1}},
\qquad -1\leq\delta\leq1.
\end{equation}
بنابراین مختصات نهایی فیدوشال با دقت زیرپیکسلی گزارش می‌شود
\cite{debella2011subpixel}.

\subsection{مرحله پنجم: برآورد تبدیل‌های افاین و پروجکتیو}

\subsubsection{تبدیل افاین}

مدل افاین شش‌پارامتری به‌صورت
\begin{equation}
\begin{bmatrix}X\\Y\end{bmatrix}
=
\begin{bmatrix}
a_1&a_2&a_3\\
b_1&b_2&b_3
\end{bmatrix}
\begin{bmatrix}r\\c\\1\end{bmatrix}
\end{equation}
تعریف و با استفاده از هر هشت فیدوشال در چارچوب کمترین مربعات حل می‌شود. این مدل
انتقال، دوران، مقیاس‌های مستقل و برش را پوشش می‌دهد
\cite{bjorck1996least,hartley2004multiple}.

\subsubsection{تبدیل پروجکتیو}

مدل پروجکتیو یا هموگرافی دارای هشت درجه آزادی است:
\begin{equation}
\lambda
\begin{bmatrix}X\\Y\\1\end{bmatrix}
=
\mathbf{H}
\begin{bmatrix}r\\c\\1\end{bmatrix}.
\end{equation}
در برنامه، مختصات هر دو دستگاه پیش از حل نرمال‌سازی می‌شوند تا وضعیت عددی دستگاه
معادلات بهتر شود. سپس ماتریس به دستگاه اصلی بازگردانده و طوری مقیاس می‌شود که
\(h_{33}=1\) باشد \cite{hartley1997eightpoint,hartley2004multiple}.

\subsubsection{معیار ارزیابی}

خطای دوبعدی   برای \(n=8\) نقطه از رابطه زیر محاسبه می‌شود:
\begin{equation}
\operatorname{RMSE}=
\sqrt{\frac{1}{n}\sum_{i=1}^{n}
\left[(X_i-\widehat X_i)^2+(Y_i-\widehat Y_i)^2\right]}.
\label{eq:rmse}
\end{equation}

\subsection{مرحله ششم: محاسبه مرکز فیدوشال}

برنامه به‌صورت خودکار نشانه‌های دارای کمینه و بیشینه \(X\) را به‌عنوان زوج افقی
\((6,8)\) و نشانه‌های دارای کمینه و بیشینه \(Y\) را به‌عنوان زوج عمودی
\((5,7)\) انتخاب می‌کند. برای هر زوج، خط ضمنی
\begin{equation}
\alpha r+\beta c+\gamma=0
\end{equation}
از دو نقطه عبور داده می‌شود. حل دستگاه دو خط، مختصات مرکز فیدوشال را در تصویر
می‌دهد. سپس اختلاف این نقطه با مرکز هندسی آرایه پیکسلی محاسبه و با تبدیل افاین به
میلی‌متر تبدیل می‌شود.

\section{نتایج عددی}

\subsection{پارامترهای تبدیل}

پارامترهای تبدیل تشابهی اولیه به‌صورت زیر گزارش شده‌اند:
\begin{equation}
[a,b,t_x,t_y]=
[-3.7518346050{\times}10^{-4},
-1.9995770203{\times}10^{-2},
-120.46154830,
126.70923615].
\end{equation}

ماتریس افاین نهایی برابر است با:
\begin{equation}
\mathbf{A}=
\begin{bmatrix}
-3.7674850528{\times}10^{-4} &
 1.9997776779{\times}10^{-2} &
-120.45964375\\
-1.9996938147{\times}10^{-2} &
-3.7496573057{\times}10^{-4} &
126.71769865
\end{bmatrix}.
\end{equation}

هموگرافی پروجکتیو نیز به‌صورت زیر به‌دست آمده است:
\begin{equation}
\mathbf{H}=
\begin{bmatrix}
-3.7673823807{\times}10^{-4} &
 1.9997231059{\times}10^{-2} &
-120.45874660\\
-1.9996392498{\times}10^{-2} &
-3.7495550541{\times}10^{-4} &
126.71383007\\
 1.0548191244{\times}10^{-9} &
-5.5122046814{\times}10^{-9} &
1
\end{bmatrix}.
\end{equation}

اندازه بردارهای ستونی بخش خطی ماتریس افاین نشان می‌دهد اندازه پیکسل در دو راستا
تقریباً برابر با \(0.02000~\unit{mm/px}\)، یعنی حدود \(20~\unit{\mu m}\) است.
بر این اساس، گستره تقریبی تصویر روی صفحه فیلم حدود
\(243.47\times244.08\,\mathrm{mm}\) و پوشش زمینی اسمی آن در مقیاس \(1:5000\) نزدیک
به \(1.217\times1.220~\unit{km}\) است.

\subsection{مختصات آشکارشده و کیفیت تطبیق}

\begin{table}[H]
\centering
\caption{مختصات پیکسلی نهایی فیدوشال‌ها و امتیاز همبستگی}
\label{tab:detections}
\begin{tabular}{cccc}
\toprule
\textbf{شناسه} & \textbf{سطر [پیکسل]} & \textbf{ستون [پیکسل]} & \lr{\textbf{NCC}} \\
\midrule
1 & \(11421.4498\) & \(11539.8249\) & \(0.812665\) \\
2 & \(11620.6411\) & \(942.0369\)   & \(0.817536\) \\
3 & \(1021.6030\)  & \(742.2988\)   & \(0.800569\) \\
4 & \(823.6901\)   & \(11340.2431\) & \(0.805463\) \\
5 & \(11719.9596\) & \(6244.2469\)  & \(0.853976\) \\
6 & \(6325.0473\)  & \(642.4465\)   & \(0.868378\) \\
7 & \(722.5797\)   & \(6037.1893\)  & \(0.877461\) \\
8 & \(6118.3165\)  & \(11639.5328\) & \(0.879561\) \\
\bottomrule
\end{tabular}
\end{table}

تمام امتیازها بین \(0.8006\) و \(0.8796\) و بسیار بیشتر از آستانه پذیرش \(0.35\)
هستند. میانگین امتیازها حدود \(0.8395\) است؛ بنابراین تطبیق‌ها از نظر معیار داخلی
برنامه قابل اعتماد ارزیابی می‌شوند. البته مقدار بالای \lr{NCC} به‌تنهایی جایگزین
کنترل بصری یا بررسی باقی‌مانده مختصات هر نقطه نیست.

\subsection{مقایسه دقت مدل‌ها}

\begin{table}[H]
\centering
\caption{مقایسه خطای تبدیل‌ها}
\label{tab:rmse}
\begin{tabular}{ccc}
\toprule
\textbf{مدل} & \(\operatorname{RMSE}~[\unit{mm}]\) &
\textbf{معادل تقریبی [پیکسل]} \\
\midrule
افاین & \(0.00972467\) & \(0.486\) \\
پروجکتیو & \(0.00936102\) & \(0.468\) \\
\bottomrule
\end{tabular}
\end{table}

مدل پروجکتیو نسبت به مدل افاین حدود \(3.74\) درصد کاهش \lr{RMSE} ایجاد کرده است.
با وجود این، عناصر پروجکتیو سطر سوم ماتریس \(\mathbf H\) در مرتبه
\(10^{-9}\) هستند؛ بنابراین مؤلفه پرسپکتیوی بسیار کوچک و هندسه تصویر به مدل افاین
نزدیک است. از دید کاربردی، هر دو مدل دقتی بهتر از نیم پیکسل ارائه می‌کنند. برتری
اندک مدل پروجکتیو باید با احتیاط تفسیر شود، زیرا این مدل با هشت پارامتر و هشت نقطه،
آزادی بیشتری برای جذب خطاهای اندازه‌گیری دارد.

\subsection{مرکز فیدوشال و جابه‌جایی آن}

نتیجه تقاطع زوج‌های روبه‌رو به‌صورت زیر است:
\begin{align}
(r_0,c_0)&=(6221.686846,\ 6140.725930)~\unit{px},\\
(\Delta r,\Delta c)&=(135.186846,\ 39.225930)~\unit{px},\\
(\Delta X,\Delta Y)&=(0.73349996,\ -2.71803138)\,\mathrm{mm}.
\end{align}
مرکز هندسی آرایه تصویر در قرارداد برنامه برابر
\((6086.5,\ 6101.5)\) پیکسل است. بنابراین مرکز فیدوشال نسبت به مرکز آرایه اسکن‌شده
حدود \(135.19\) پیکسل در راستای افزایش سطر و \(39.23\) پیکسل در راستای افزایش ستون
جابجا است. این مقدار بیشتر بیانگر جابه‌جایی قاب فیلم در فرایند اسکن و برش تصویر است
و نباید مستقیماً با آفست بسیار کوچک \lr{PPS} در گواهی کالیبراسیون یکسان تلقی شود.

\begin{figure}[H]
\centering
\includegraphics[width=1\textwidth]{Figure_1.png}
\caption{مکان هشت فیدوشال آشکارشده و مرکز فیدوشال محاسبه‌شده روی تصویر هوایی}
\label{fig:detected-image}
\end{figure}

\section{بحث و ارزیابی}

\subsection{نقاط قوت روش}

\begin{itemize}[rightmargin=1.2em]
  \item تنها دو اندازه‌گیری اولیه از کاربر گرفته می‌شود و ادامه فرایند خودکار است.
  \item جست‌وجوی محلی پیرامون مکان پیش‌بینی‌شده، احتمال تطبیق اشتباه با عوارض زمینی
  را کاهش می‌دهد.
  \item پالایش سهموی قله همبستگی، مختصات زیرپیکسلی بدون بهینه‌سازی سنگین فراهم می‌کند.
  \item استفاده هم‌زمان از هر هشت فیدوشال، افزونگی لازم برای برآورد افاین و ارزیابی
  \lr{RMSE} را ایجاد می‌کند.
  \item نرمال‌سازی مختصات پیش از حل هموگرافی، پایداری عددی برآورد پروجکتیو را بهبود
  می‌دهد.
  \item کنترل امتیاز حداقل و ممنوعیت پذیرش قله روی مرز پنجره، دو آزمون ساده اما مؤثر
  برای شناسایی شکست تطبیق هستند.
\end{itemize}

\subsection{محدودیت‌ها و ملاحظات تفسیری}

\begin{enumerate}[rightmargin=1.4em]
  \item \textbf{عدم اعمال کامل گواهی کالیبراسیون:}
  تابع خواندن \code{RC20.txt} در نسخه فعلی فقط بخش \code{$FIDUCIALS} را وارد
  محاسبات می‌کند. فاصله کانونی، \lr{PPS} و جدول اعوجاج شعاعی خوانده یا اعمال
  نمی‌شوند. بنابراین خروجی فعلی یک تبدیل فیدوشال‌محور است، نه یک مدل کامل
  کالیبراسیون داخلی برای تشکیل مختصات تصحیح‌شده پرتو
  \cite{brown1971calibration,wolf2014elements}.

  \item \textbf{معنای نقطه گزارش‌شده:}
  نقطه‌ای که برنامه با عنوان «نقطه اصلی تصحیح‌شده» نمایش می‌دهد، از نظر هندسی
  \emph{مرکز فیدوشال} حاصل از تقاطع خطوط \(6\!-\!8\) و \(5\!-\!7\) است. برای
  دستیابی به \lr{PPS} کالیبره باید بردار
  \((-0.004,+0.005)\,\mathrm{mm}\) نیز با قرارداد محورها و علامت‌های صحیح به آن افزوده
  شود. این تصحیح فقط حدود یک‌چهارم پیکسل است، اما از نظر مفهومی مهم است.

  \item \textbf{الگوی مصنوعی:}
  تطبیق با حلقه مصنوعی نسبت به تغییر شکل، نوردهی و آلودگی نشانه حساس است. استفاده از
  الگوهای مستقل و واقعی برای هر نوع فیدوشال می‌تواند استحکام روش را افزایش دهد.

  \item \textbf{آمار دقت محدود:}
  برنامه \lr{RMSE} کلی را گزارش می‌کند، اما باقی‌مانده دوبعدی هر فیدوشال، انحراف
  معیار پارامترها و آزمون نقاط پرت را ارائه نمی‌دهد. افزودن این موارد برای یک
  ارزیابی فتوگرامتری کامل توصیه می‌شود.

  \item \textbf{مقایسه افاین و پروجکتیو:}
  کاهش کوچک خطای مدل پروجکتیو الزاماً به معنای برتری فیزیکی آن نیست؛ زیرا تعداد
  پارامترهای بیشتر می‌تواند بخشی از نویز اندازه‌گیری را جذب کند. برای انتخاب مدل،
  علاوه بر \lr{RMSE} باید الگوی باقی‌مانده‌ها و منطق هندسی اسکنر نیز بررسی شود.
\end{enumerate}

\section{پیشنهاد برای تکمیل مدل توجیه داخلی}

برای تبدیل این برنامه به یک سامانه کامل‌تر، مراحل زیر پیشنهاد می‌شوند:
\begin{enumerate}[rightmargin=1.4em]
  \item خواندن فاصله کانونی، \lr{PPS} و جدول اعوجاج از \code{RC20.txt}؛
  \item تبدیل مختصات پیکسلی همه نقاط تصویر به مختصات متریکی با مدل منتخب؛
  \item انتقال مبدأ از مرکز فیدوشال به نقطه اصلی کالیبره؛
  \item درون‌یابی جدول اعوجاج شعاعی و تصحیح مختصات بر حسب شعاع
  \cite{brown1971calibration}؛
  \item گزارش بردار باقی‌مانده هر فیدوشال و ترسیم آن با بزرگ‌نمایی مناسب؛
  \item استفاده از برآورد مقاوم یا کنترل خودکار نقاط پرت؛
  \item مقایسه مدل‌ها با معیارهای سرشکن‌شده بر حسب تعداد پارامترها یا اعتبارسنجی
  متقابل.
\end{enumerate}

\section{جمع‌بندی}

در این پروژه، یک روش نیمه‌خودکار و منظم برای شناسایی هشت فیدوشال مارک و برقراری
رابطه میان مختصات پیکسلی و مختصات متریک دوربین اجرا شد. دو فیدوشال اولیه برای تعیین
تقریب هندسی استفاده شدند و سایر فیدوشال‌ها با تطبیق \lr{NCC} و پالایش زیرپیکسلی
اندازه‌گیری شدند. همه امتیازهای همبستگی بیشتر از \(0.80\) بودند. مدل افاین خطای
\(0.00972\,\mathrm{mm}\) و مدل پروجکتیو خطای \(0.00936\,\mathrm{mm}\) ایجاد کرد؛ هر دو
مقدار تقریباً معادل نیم پیکسل هستند. مرکز فیدوشال نیز در مختصات
\((6221.687,\ 6140.726)\) پیکسل تعیین شد.

نتایج نشان می‌دهند گردش‌کار برای بازیابی هندسه قاب و اندازه‌گیری خودکار فیدوشال‌ها
موفق است. با این حال، برای آن‌که خروجی به‌عنوان توجیه داخلی کامل دوربین به کار رود،
لازم است آفست \lr{PPS} و تصحیحات اعوجاج شعاعی ثبت‌شده در گواهی کالیبراسیون نیز به
مدل افزوده شوند.

\section*{منابع داده و نرم‌افزار}
\addcontentsline{toc}{section}{منابع داده و نرم‌افزار}

\begin{itemize}[rightmargin=1.2em]
  \item \code{interior_orientation.py}: پیاده‌سازی مراحل پردازش و محاسبات.
  \item \code{RC20.txt}: مشخصات دوربین، مختصات فیدوشال‌ها و جدول اعوجاج.
  \item \code{meta.txt}: مشخصات پرواز، مقیاس عکس و ارتفاع منطقه.
  \item \code{results.txt}: پارامترهای تبدیل، مختصات آشکارشده و نتایج دقت.
  \item \code{1614.tif}: تصویر هوایی اسکن‌شده شهر همدان.
\end{itemize}
\newpage
\printbibliography
\end{document}
